\subsection{P-16. An explicit width-three obstruction to uncentered two-pair interpolation}
\label{P-16}

\paragraph{Statement and scope.}
Work in characteristic zero. There are eighteen explicitly specified,
pairwise distinct algebraic Lagrange nodes and a nonzero multilinear
polynomial $H$ of degree $16$, of minimum ROABP width exactly three,
such that
\[
 H\bigl(\operatorname{SV}_2(\mathbf y,\mathbf z)\bigr)=0.
 \tag{P16.1}
\]
Its displayed width-three representation has every univariate factor
generically invertible. In contrast, P-14 proves
\[
 H\bigl(t\mathbf1+\operatorname{SV}_2(\mathbf y,\mathbf z)\bigr)\ne0.
 \tag{P16.2}
\]
Thus this is an obstruction to the uncentered map, not to balanced
$\operatorname{RFE}_{2}^{2}$ or to centered two-pair interpolation.
The latter two maps on the full width-three class with rank-two resets
remain separate questions.

We use $\operatorname{RFE}_{k}^{\ell}$ for evaluations of rational
functions whose numerator degree is at most $k$ and denominator degree
at most $\ell$. In fact the construction below gives, for every
$k\ge1$, a quadratic width-three annihilator of
$\operatorname{RFE}_{k}^{1}$ on $6(k+1)$ nodes, and its multilinear
reversal annihilates $\operatorname{RFE}_{1}^{k}$. Only $k=2$ is needed
for (P16.1).

\paragraph{Explicit disjoint pole groups.}
Put $d=k+1$, choose positive real radicals
\[
 \alpha=2^{1/(2d)},\qquad \beta=(3/2)^{1/(2d)},
 \qquad a=\beta^d=\sqrt{3/2},
\]
and let $\zeta$ be a primitive $2d$th root of unity. Work over the
number field
\[
 E=\mathbb Q(\zeta,\alpha,\beta),\qquad
 [E:\mathbb Q]\le(2d)^2\varphi(2d).
 \tag{P16.3}
\]
This bound follows by adjoining the two radicals, each of degree at
most $2d$, to the cyclotomic field. Define
\[
 \begin{array}{ll}
 P_1(T)=T^{2d}-1,&P_2(T)=T^d-a,\\
 P_3(T)=-2(T^d+a),&P_4(T)=T^{2d}-2.
 \end{array}
\]
Let $\mathcal A_i$ be the roots of $P_i$. Explicitly these four
ordered groups can be listed as
\[
 \begin{split}
 \mathcal A_1&=(\zeta^j)_{0\le j<2d},\qquad
 \mathcal A_2=(\beta\zeta^{2j})_{0\le j<d},\\
 \mathcal A_3&=(\beta\zeta^{2j+1})_{0\le j<d},\qquad
 \mathcal A_4=(\alpha\zeta^j)_{0\le j<2d}.
 \end{split}
 \tag{P16.4}
\]
All roots are simple. The first, middle, and last groups have distinct
absolute values $1,\beta,\alpha$, and the two middle root sets are
separated by their $d$th powers $a$ and $-a$. Thus all $6d$ nodes
are distinct. For $k=2$, $d=3$, their field has degree at most $72$.

Set
\[
 R_1=\frac{P_2}{P_1},\qquad R_2=\frac1{P_2},\qquad
 R_3=\frac1{P_3},\qquad R_4=\frac{P_3}{P_4}.
 \tag{P16.5}
\]
Each $R_i$ has precisely the simple poles in $\mathcal A_i$, with
nonzero residues, and satisfies $R_i(T)=O(T^{-d})$ at infinity.
Moreover $P_2P_3=-P_1-P_4$, and hence
\[
 R_1R_2+R_1R_4+R_3R_4
 =\frac{P_4+P_2P_3+P_1}{P_1P_4}=0.
 \tag{P16.6}
\]
For $b\in\mathcal A_i$, define the explicit nonzero residue
$\rho_b=\lim_{T\to b}(T-b)R_i(T)$. Introduce one variable $X_b$ for
each node and put
\[
 L_i(\mathbf X)=\sum_{b\in\mathcal A_i}\rho_bX_b,
 \qquad
 Q=L_1L_2+L_1L_4+L_3L_4.
 \tag{P16.7}
\]
The four variable sets are disjoint, so $Q$ is nonzero, multilinear,
and homogeneous of degree two.

\paragraph{The unbalanced rational-evaluation identity.}
The partial fraction identity
\[
 R_i(T)=\sum_{b\in\mathcal A_i}\frac{\rho_b}{T-b}
\]
and its expansion at infinity give
\[
 \sum_{b\in\mathcal A_i}\rho_b b^j=0
       \quad(0\le j\le d-2=k-1).
 \tag{P16.8}
\]
A generic rational function with numerator degree at most $k$ and
linear denominator has the form
\[
 r(T)=p(T)+\frac{c}{T-u},\qquad \deg p\le k-1.
\]
On the open set on which no evaluation denominator vanishes,
(P16.8) implies
\[
 L_i\bigl((r(b))_b\bigr)=-cR_i(u).
\]
Equation (P16.6) proves that $Q$ vanishes on these evaluations.
Clearing the finitely many evaluation denominators gives a polynomial
identity in the numerator and denominator coefficients. Its vanishing
on the dense open set of genuine linear denominators therefore extends
to all defined evaluations in $\operatorname{RFE}_{k}^{1}$, including
constant denominators. No assertion that a degree-$k$ polynomial part
is killed term by term is needed for this last specialization.

\paragraph{The exact width and its reversal.}
Read the variables in the four group orders in (P16.4). For each node,
use
\[
 M_b(X_b)=
 \begin{cases}
 I_3+\rho_bX_bE_{12},&b\in\mathcal A_1\cup\mathcal A_3,\\
 I_3+\rho_bX_bE_{23},&b\in\mathcal A_2\cup\mathcal A_4.
 \end{cases}
 \tag{P16.9}
\]
The $(1,3)$ entry of their product is exactly $Q$: the possible
ordered $E_{12}E_{23}$ pairs come from groups $(1,2),(1,4),(3,4)$,
and products of three such matrix units are zero. Thus $Q$ has a
width-three unit-factor representation. Its minimum width in any
order cannot be two: it is a nonzero annihilator of balanced
$\operatorname{RFE}_{1}^{1}$, whereas P-14 hits every nonzero
width-two ROABP with that family. Alternatively, its rank in the
specified order is already three at a cut strictly inside group two:
the independent left polynomials $1,L_1,L_1L_{2,\mathrm{left}}$
are paired with independent right polynomials
$L_3L_4,L_{2,\mathrm{right}}+L_4,1$, of three different degrees.

Let $N=6d$, and reverse the coefficient basis in every variable:
\[
 Q^{\mathrm{rev}}(\mathbf X)=
       \left(\prod_bX_b\right)Q\bigl((X_b^{-1})_b\bigr).
 \tag{P16.10}
\]
It is a nonzero multilinear polynomial of degree $N-2$. For every
variable partition, reversal merely permutes the rows and columns
of the coefficient matrix. Hence it preserves every coefficient-cut
rank and the minimum width in every order. Replacing (P16.9) by
$X_bI_3+\rho_bE_{12}$ or $X_bI_3+\rho_bE_{23}$ gives its explicit
width-three word; each factor has determinant $X_b^3$.

If $r$ has numerator degree at most one and denominator degree at
most $k$, then $1/r$ has numerator degree at most $k$ and denominator
degree at most one. On the dense open set where all relevant values
are nonzero and defined, (P16.10) therefore vanishes on $(r(b))_b$.
Clearing denominators extends the identity to every defined
$\operatorname{RFE}_{1}^{k}$ evaluation.

\paragraph{The ordinary SV convention and the repair by centering.}
Specialize now to $k=2$, $d=3$, $N=18$. Let
\[
 A(T)=\prod_b(T-b),\qquad D_b=A'(b)\ne0,
 \qquad \lambda_b(Y)=\frac{A(Y)}{(Y-b)D_b},
 \qquad H(\mathbf X)=Q^{\mathrm{rev}}\bigl((D_bX_b)_b\bigr).
 \tag{P16.11}
\]
For the ordinary two-pair generator
$\operatorname{SV}_{2,b}=z_1\lambda_b(y_1)+z_2\lambda_b(y_2)$, one has
\[
 D_b\operatorname{SV}_{2,b}
   =-\sum_{j=1}^{2}\frac{z_jA(y_j)}{b-y_j}.
 \tag{P16.12}
\]
As a function of the node $b$, the right-hand side is a rational
function with numerator degree at most one and denominator degree
at most two. The proved $\operatorname{RFE}_{1}^{2}$ identity gives
(P16.1) on a dense open parameter set and thus identically. This
also agrees with the fixed diagonal equivalence in Hu et al. (2024),
Proposition~2.2, pp.~16--17 (P-B20).

The displayed factors for $H$ are
$D_bX_bI_3+\rho_bE_{12}$ or $D_bX_bI_3+\rho_bE_{23}$, with determinant
$D_b^3X_b^3$. All are generically invertible, and diagonal variable
scaling preserves all coefficient-cut ranks. Thus no variable order permits width two for $H$, while the displayed
order attains width three. P-14, with $w=3$ and center slopes
all one, applies directly to this word and proves (P16.2).

This construction does not invoke a monotonicity assertion for the
width of factors. Its width claims follow from explicit matrix words
and exact coefficient-basis permutations. It also does not establish
failure of the centered map on rank-two reset words; that is the
remaining balanced question.
\hfill$\square$
